\documentclass[11pt]{article}
\newcommand{\ShortTitle}{LEHMER / RAMANUJAN TAU}
\usepackage[T1]{fontenc}
\usepackage[utf8]{inputenc}
\usepackage{mathpazo}
\usepackage[scaled=0.92]{helvet}
\usepackage{courier}
\usepackage{amsmath,amssymb,mathtools}
\usepackage[a4paper,landscape,left=18mm,right=18mm,top=19mm,bottom=17mm,headheight=14pt]{geometry}
\usepackage{graphicx,microtype,xcolor,fancyhdr,hyperref,booktabs,array}
\definecolor{Ink}{HTML}{243138}
\definecolor{Teal}{HTML}{21676B}
\definecolor{Copper}{HTML}{A45C36}
\definecolor{Soft}{HTML}{E9F0EF}
\definecolor{Gray}{HTML}{627078}
\hypersetup{colorlinks=true,linkcolor=Teal,citecolor=Teal,urlcolor=Teal,pdfauthor={Komorebi P.}}
\urlstyle{same}
\setlength{\parindent}{0pt}
\setlength{\parskip}{0.55em}
\setlength{\emergencystretch}{1.5em}
\linespread{1.04}
\pagestyle{fancy}
\fancyhf{}
\fancyhead[L]{\sffamily\scriptsize\color{Gray}OPERATION KOMOREBI / VISUAL GUIDE}
\fancyhead[R]{\sffamily\scriptsize\color{Gray}\ShortTitle}
\fancyfoot[L]{\sffamily\scriptsize\color{Gray}READER EDITION 1.0 / 9 OCTOBER 2026}
\fancyfoot[R]{\sffamily\small\color{Ink}\thepage}
\renewcommand{\headrulewidth}{0.25pt}
\renewcommand{\footrulewidth}{0pt}
\providecommand{\N}{\mathbb N}
\providecommand{\Z}{\mathbb Z}
\providecommand{\Q}{\mathbb Q}
\providecommand{\R}{\mathbb R}
\providecommand{\C}{\mathbb C}
\providecommand{\eps}{\varepsilon}
\providecommand{\dd}{\mathrm d}
\providecommand{\norm}[1]{\left\lVert #1\right\rVert}
\providecommand{\abs}[1]{\left\lvert #1\right\rvert}
\newcommand{\GuideCover}[4]{%
\hypersetup{pageanchor=false}
\begin{titlepage}\thispagestyle{empty}
{\sffamily\small\color{Teal}OPERATION KOMOREBI\hfill CONJECTURE SERIES #1 / VISUAL GUIDE}\par
\vspace{18mm}
{\fontsize{40}{44}\selectfont\color{Ink}#2}\par
\vspace{4mm}{\LARGE\color{Ink}#3}\par
\vspace{6mm}{\color{Teal}\rule{\linewidth}{1.1pt}}\par
\vspace{4mm}{\large Komorebi P.}\par
{\sffamily\small\color{Gray}READER EDITION 1.0 / 9 OCTOBER 2026}\par
\vfill
\begin{minipage}{0.78\linewidth}\large #4\end{minipage}\par
\vspace{5mm}{\sffamily\small\color{Teal}WORKED EXAMPLES\quad /\quad EXACT COORDINATES\quad /\quad OPEN PROBLEMS}\par
\end{titlepage}\hypersetup{pageanchor=true}}
\newcommand{\GuideSpread}[5]{%
\clearpage
{\fontsize{22}{26}\selectfont\bfseries\color{Ink}#1}\par
\vspace{1mm}{\sffamily\small\color{Gray}#2}\par
\vspace{3mm}
\noindent\begin{minipage}[t]{0.665\linewidth}\vspace{0pt}
\centering\includegraphics[width=\linewidth,height=113mm,keepaspectratio]{#3}
\end{minipage}\hfill
\begin{minipage}[t]{0.295\linewidth}\vspace{0pt}
{\sffamily\small\bfseries\color{Teal}READ THE FIGURE}\par
\small #4
\end{minipage}\par
\vfill
\noindent\colorbox{Soft}{\parbox{\dimexpr\linewidth-2\fboxsep\relax}{\vspace{1.3mm}\small\textbf{Carry forward.} #5\vspace{1.3mm}}}\par
}

\hypersetup{pdftitle={Lehmer's Conjecture: A Visual Guide},pdfsubject={Operation Komorebi, Conjecture Series 07}}
\newcommand{\V}{\mathcal V}
\newcommand{\ch}{\chi_5}
\begin{document}
\GuideCover{07}{Lehmer's conjecture}{Ramanujan tau, signed packets, and the prime gap}{Thirty visual studies of one fixed modular form. Exact coefficients, classical identities, project arguments and constructed local models are distinguished throughout. The program has not proved that every coefficient is nonzero.}

\GuideSpread{One product, infinitely many coefficients}{DEFINITION / THE FIXED MODULAR FORM}{../shared/figures/tau_product.pdf}{
The coefficient of $q^n$ is $\tau(n)$. The initial $q$ shifts the product coefficient by one: $\tau(n)=[q^{n-1}]\prod_{m\ge1}(1-q^m)^{24}$.

Lehmer's conjecture asks whether $\tau(n)\ne0$ for every positive integer $n$. Large values and frequent sign changes do not answer that question.

The target is the actual normalized weight-12 cusp form $\Delta$ of level one.}{A sequence constructed to obey some of the same rules is a separate mathematical object.}

\GuideSpread{Three independent coefficient calculations}{EXACT FINITE EVIDENCE / ALL INTEGERS THROUGH 512}{../shared/figures/coefficient_methods.pdf}{
The replay compares direct product expansion, a Jacobi triangular recurrence, and a divisor-sum recurrence from the logarithmic derivative.

If $F=\Delta/q=\sum a_mq^m$, the last method uses
\[
m a_m=-24\sum_{j=1}^{m}\sigma_1(j)a_{m-j}.
\]
All three integer methods agree on $512$ coefficients. Every one is nonzero.}{Agreement validates this bounded table. It is not a proof beyond the tested range.}

\GuideSpread{The actual coefficients through 512}{EXACT INPUTS / SIGNED LOGARITHMIC DISPLAY}{../shared/figures/coefficients512.pdf}{
The vertical coordinate compresses large integers while preserving their sign. Teal points are positive and copper points negative.

Every input is an exact integer. Floating-point logarithms display the magnitudes only.

This plot includes composite indices and primes. It is not a normalized prime distribution or a claimed limiting law.}{A scatter plot supports exploration and checking. It cannot resolve the universal nonzero claim.}

\GuideSpread{Deligne's bound controls size}{CLASSICAL THEOREM / PRIME DATA THROUGH 512}{../shared/figures/prime_bound.pdf}{
For every prime $p$, Deligne gives
\[
|\tau(p)|\le2p^{11/2}.
\]
The plot divides each prime coefficient by this upper bound. All displayed values lie between the two horizontal lines.

Zero remains inside the allowed interval. The bound is fully compatible with a hypothetical zero coefficient.}{The remaining step requires a nonzero conclusion. An upper bound supplies no positive lower bound.}

\GuideSpread{Multiplicativity organizes the indices}{CLASSICAL IDENTITY / AN EXACT WORKED EXAMPLE}{../shared/figures/multiplicativity.pdf}{
If $\gcd(m,n)=1$, then $\tau(mn)=\tau(m)\tau(n)$. Here $4$ and $3$ are coprime and their exact coefficients give $\tau(12)=-370944$.

Unique prime factorization reduces all coefficients to prime-power blocks.

Coprimality matters: $\tau(4)$ is not $\tau(2)^2$. It uses the Hecke recurrence.}{Multiplicativity identifies the correct local blocks. It does not establish their nonvanishing.}

\GuideSpread{The recurrence lives at one prime}{CLASSICAL HECKE IDENTITY / NOT INTERPOLATION IN p}{../shared/figures/hecke_recurrence.pdf}{
Fix a prime $p$. Write $T_r=\tau(p^r)$ and $a=\tau(p)$. The recurrence starts from $T_0=1$, $T_1=a$ and applies for $r\ge2$.

The roots of $X^2-aX+p^{11}$ encode the same sequence.

This varies the exponent $r$. It supplies no single analytic function of the prime index $p$.}{Interpolation across primes needs its own definition and proof. The exponent recurrence does not supply it.}

\GuideSpread{Why settling primes is enough}{CLASSICAL REDUCTION / THE ROOT-RATIO ARGUMENT}{../shared/figures/prime_reduction.pdf}{
A prime-power zero forces the Hecke-root ratio to be a root of unity. The displayed rational algebraic integer can only be $0,1,2,3,4$.

For $p\ge5$, each nonzero alternative would make an integer square have odd $p$-adic valuation. The known values at $2$ and $3$ exclude their alternatives directly.

Multiplicativity completes the reduction.}{The remaining target is $\tau(p)\ne0$ at every prime, with no unexamined prime-power exception.}

\GuideSpread{Parity has a precise square pattern}{CLASSICAL IDENTITY / EXACT EXAMPLES THROUGH 81}{../shared/figures/parity.pdf}{
Jacobi's identity yields
\[
\tau(n)\text{ odd}\iff n\text{ an odd square}.
\]
The marked indices are $1,9,25,49,81$. Every other displayed coefficient is even.

In particular an odd prime has an even coefficient. Evenness does not decide whether that integer is zero.}{Parity rules out particular values. It leaves the prime zero value permitted.}

\GuideSpread{Hook lengths select the five-cores}{CLASSICAL HOOK IDENTITY / TWO PARTITIONS OF 6}{../shared/figures/fivecore_hooks.pdf}{
Han's hook formula uses $\prod_h(1-z/h^2)$. At $z=25$, a hook of length $5$ kills the contribution.

The surviving partitions are the 5-cores. The hooks of $(4,1,1)$ are $6,3,2,1,2,1$, so it survives. The row $(6)$ contains a hook of length $5$.

Partitions of size $n-1$ contribute to $\tau(n)$.}{The formula exactly re-expresses the coefficient. Surviving weights still have mixed signs.}

\GuideSpread{A packet is a five-coordinate lattice point}{CLASSICAL FIVE-CORE CODING / THE UNIQUE PACKET AT n=2}{../shared/figures/shell_coordinates.pdf}{
The shell is
\[
\begin{gathered}
\V_n=\{v\in\Z^5:\ v_i\equiv i\pmod5,\\
\sum v_i=0,\quad \sum v_i^2=10n\}.
\end{gathered}
\]
Labels run from $0$ to $4$. Han's normalization is $|\lambda|=\|v\|^2/10-1$.

The unique packet at $n=2$ has weight $-24$, exactly $\tau(2)$.}{The shell is finite. A projection would not represent its full four-dimensional geometry.}

\GuideSpread{Vandermonde weights are integers}{EXACT ALGEBRA / EVERY PACKET WEIGHT IS NONZERO}{../shared/figures/vandermonde.pdf}{
There are ten differences and denominator $288$. Reversing all ten differences preserves the sign.

The binomial determinant proves integrality, including for negative coordinates. Distinct residues modulo $5$ make all differences nonzero.

Every packet weight is $1$ modulo $5$. Thus each is a nonzero integer, even when their signed total might vanish.}{Nonzero summands can cancel. The question concerns their total.}

\GuideSpread{At 7, three orbit masses compete}{EXACT SHELL ENUMERATION / ORIGINAL MULTIPLICITIES RETAINED}{../shared/figures/orbit7.pdf}{
The involution
\[
J(v)_i=-v_{(-i\bmod5)}
\]
preserves the shell and $W$. A bar sums one entire $J$-orbit.

The six packets at $7$ form three orbits with mixed signs. Their exact total is $-16744$.

Pairing by $J$ simplifies bookkeeping. It does not match positive weights with negative weights.}{An orbit weight includes one or two original packets. Retain its multiplicity in every total.}

\GuideSpread{At 31, mixed signs remain visible}{EXACT SHELL ENUMERATION / 16 J-ORBITS}{../shared/figures/orbit31.pdf}{
The prime shell at $31$ has $32$ packets and $16$ orbits. Eight orbit masses are positive and eight negative.

Their total is
\[
\tau(31)=-52843168.
\]
This is an exact finite coefficient calculation. Equal counts of positive and negative terms do not imply cancellation.

The bar height is mass divided by $1000$.}{Resolving all orbits exposes the arithmetic. It does not automatically choose a dominant sign or valuation.}

\GuideSpread{Packet counts recover a congruence}{CLASSICAL COUNT / EXACT PRIME FORMULA}{../shared/figures/packet_count.pdf}{
The five-core count is
\[
|\V_n|=\sum_{d\mid n}\ch(d)\frac nd,
\]
where $\ch(d)=(d/5)$, zero on multiples of $5$.

At a prime it equals $p+\ch(p)$. Since every weight is $1$ modulo $5$, the weighted sum has the same residue.

For $p\ne5$, a zero requires $p\equiv4\pmod5$.}{The congruence excludes three classes. An infinite residue class still needs additional arithmetic.}

\GuideSpread{Content creates a conserved split}{PROJECT ARGUMENT / EXACT MASSES AT n=121}{../shared/figures/content121.pdf}{
Content is the positive gcd $g(v)=\gcd(v_0,\ldots,v_4)$. Its square divides $n$, and $5$ does not divide it.

At $121$, content-$1$ and content-$11$ masses have opposite signs. Their total is $498319933$.

Every finite symmetric polynomial correction that is universally null conserves each content mass. It cannot make all corrected masses share one sign here.}{This obstruction concerns multiple content classes. A prime shell has only content $1$.}

\GuideSpread{Primitive extraction has an exact inverse}{PROJECT ALGEBRA / MOBIUS INVERSION ON SQUARE DIVISORS}{../shared/figures/mobius_residual.pdf}{
The primitive sum is
\[
R(n)=\sum_{v\in\V_n,\ g(v)=1}W(v).
\]
Rescaling an exact-content packet multiplies its weight by $\ch(g)g^{10}$. The displayed formulas are inverse convolutions.

The convention $\ch(g)=0$ on multiples of $5$ avoids nonexistent content sectors. At a prime, $R(p)=\tau(p)$.}{Primitive extraction classifies radial contributions. It preserves the original prime obstruction.}

\GuideSpread{The primitive local generating factor}{PROJECT ALGEBRA / FORMAL POWER SERIES AT A FIXED PRIME}{../shared/figures/local_factor.pdf}{
The denominator is the Hecke denominator. The numerator removes packets scaled by $p$.

The first coefficients are $R(1)=1$ and $R(p)=\tau(p)$. The subtraction formula starts at $r=2$.

The denominator is invertible as a formal series because its constant coefficient is $1$. No convergence assumption is needed.}{Study the fixed-prime recurrence without assuming that its first nonconstant coefficient is nonzero.}

\GuideSpread{Squares are settled for the primitive residual}{PROJECT THEOREM / FINITE ILLUSTRATION AT s=1 THROUGH 22}{../shared/figures/square_residuals.pdf}{
The project proves $R(p^{2m})\ne0$ for every prime $p$ and $m\ge1$. Multiplicativity gives $R(s^2)\ne0$ for every $s\ge1$.

The proof includes $a=\tau(p)=0$. For $p\equiv4\pmod5$, it then gives
\[
R(p^{2m})=(-1)^mp^{11m-1}(p+1),
\]
which is nonzero.

The plotted values illustrate the written proof.}{The square theorem explicitly permits a zero prime parameter. It cannot force that parameter to be nonzero.}

\GuideSpread{A null transfer preserves the signed total}{CONSTRUCTED LINEAR-ALGEBRA EXAMPLE / NOT PACKET DATA}{../shared/figures/null_transfer.pdf}{
The displayed masses sum to $3$. The change $(-5,+5,0)$ preserves the total and leaves three positive masses.

For actual nonzero orbit masses $w_i$, a weighted-null correction satisfies $\sum w_iq_i=0$. If $T=\sum w_i\ne0$, set
\[
q_i=T/(Nw_i)-1.
\]
Every corrected mass becomes $T/N$. This choice already uses $T$.}{Even a complete space of zero-total transfers leaves the need to prove that the total is nonzero.}

\GuideSpread{Interpolation is not universal nullness}{PROJECT ARGUMENT / THE FINITE-SHELL BOUNDARY}{../shared/figures/interpolation.pdf}{
Even symmetric coordinates identify the $J$-orbits on one shell. A rational homogeneous symmetric polynomial of degree $10(N-1)$ can interpolate arbitrary values on $N$ orbits.

That prescribes only one shell. The natural interpolated null at $7$ fails already at shell $1$.

Top-degree harmonic projection also loses values: it annihilates $s_2D$, where $D$ is the Vandermonde.}{Interpolation supplies neither universal nullness nor the missing nonzero arithmetic margin.}

\GuideSpread{Refinement separates more residue data}{SCHEMATIC TREE / c MODULO POWERS OF TWO}{../shared/figures/dyadic_tree.pdf}{
Write $c_i=(v_i-i)/5$. A depth-$k$ sector fixes $c$ modulo $2^k$, allowing $J$.

Each child belongs to one parent. The aggregation matrix $A$ gives
\[
t_k=At_{k+1},\qquad M_k=AM_{k+1}.
\]
This binary picture is schematic. Actual shells have varying occupied-child counts.}{Refinement adds constraints to sector-preserving corrections. It does not automatically add needed transfers.}

\GuideSpread{The first mixed-sign boundary is exact}{PROJECT THEOREM / ALL DYADIC DEPTHS k AT LEAST 1}{../shared/figures/mixed_threshold.pdf}{
For $M=2^k$, the first mixed shell is
\[
N_k=(20M^2+2M+\gamma_k)/7,
\]
where $\gamma_k=7,1,6$ for $k\equiv1,2,0\pmod3$.

Every shell below it has sign-pure depth-$k$ sectors. At the boundary, an explicit opposite-sign pair occupies one sector.

The logarithmic plot shows the first eight formula values.}{Sign-purity differs from orbit uniqueness. Neither statement compares global positive and negative totals.}

\GuideSpread{A global null can fail inside each sector}{EXACT COUNTEREXAMPLE / DEPTH 1 AT p=7}{../shared/figures/sector7.pdf}{
These occupied parity-sector masses are $-1848,-17248,2352$, totaling $-16744$.

The polynomial $Q_8$ is globally null. Its three sector sums are nonzero and add to zero. Multiplication by a sector indicator therefore destroys the null identity.

A correction preserving each displayed sector cannot remove their opposing signs.}{A valid localized construction must solve its parent cancellation conditions explicitly.}

\GuideSpread{A finite rank defect can disappear}{EXACT MATRIX RANKS / THE CONSTRUCTED FAMILY AT p=31}{../shared/figures/ranks31.pdf}{
There are $16$ orbit sectors and $15$ zero-total directions. The orbit-transfer ranks are $5$, $14$, $15$ at cutoffs $14$, $20$, $22$.

Degree-$22$ additions remove the remaining one-dimensional defect. The recorded rational certificate preserves coarse masses during its fine solve.

These are cutoffs in the constructed family. No optimality claim is made.}{Full rank at $31$ establishes this example. It proves neither uniform rank nor a uniformly nonzero survivor.}

\GuideSpread{Auxiliary roots use a different variable}{EXACT POLYNOMIAL EXAMPLES / n=499 AND 16249}{../shared/figures/auxiliary_roots.pdf}{
The shell polynomial has degree at most ten and $P_n(5)=288\tau(n)$. The bar in $\overline G$ denotes reduction modulo $5$.

Its integral normalization satisfies
\[
v_5(\tau(n))=\mu_n+v_5(G_n(1)).
\]
At $499$, the simple root near $1$ lifts to $16$ modulo $25$. At $16249$, every $U\equiv1\pmod5$ gives $G(U)\equiv15\pmod{25}$, so the double residual root does not lift.}{Root analysis varies $U$ and evaluates at $1$. It is not interpolation in the prime index.}

\GuideSpread{Three proposed valuation shortcuts fail}{EXACT COUNTEREXAMPLES / AUXILIARY EQUALITIES, NOT LEHMER}{../shared/figures/lifting_failures.pdf}{
The proposed identity was
\[
v_\ell(\tau(p))=v_\ell(p+1).
\]
Among primes $p\ne\ell$ with $\ell\mid(p+1)$, the first failures in the exhaustive search through $150000$ occur at $81919$, $8747$, $499$ for $\ell=2,3,5$.

The $1249$ example gives the $5$-adic discrepancy in the opposite direction. Decisive coefficients were independently checked with an Eisenstein convolution.}{These nonzero values refute proposed valuation identities. They are not counterexamples to Lehmer.}

\GuideSpread{Local filtered data allow trace zero}{CONSTRUCTED LOCAL MODEL / EVERY PRIME p}{../shared/figures/trace_phi.pdf}{
The displayed Frobenius has polynomial $X^2+p^{11}$, determinant $p^{11}$ and trace zero.

Choose a line $L$. Set $\mathrm{Fil}^iD=D$ for $i\le0$, $L$ for $1\le i\le11$, and $0$ afterward. Then $t_H=t_N=11$.

A rational stable line would require $\lambda^2=-p^{11}$ in $\Q_p$, impossible by odd valuation. The module is weakly admissible.}{The local object permits zero trace. It is not identified with the restriction of the global form $\Delta$.}

\GuideSpread{Every fixed modulus has actual survivors}{CLASSICAL-THEOREM CONSEQUENCE / ACTUAL TAU COEFFICIENTS}{../shared/figures/actual_survivors.pdf}{
Use the actual joint finite image of Deligne's integral reductions and the cyclotomic character modulo $5M$.

Complex conjugation has exact trace zero before reduction, including at auxiliary prime $2$. Chebotarev gives positive-density prime classes with the displayed residues.

The primitive factor reduces to $(1-X^2)/(1-X^2)=1$. Thus every $R(p^r)$ for $r\ge1$ is zero modulo $M$.}{No independence or image-surjectivity is assumed. A fixed list of residue tests cannot exclude all these actual primes.}

\GuideSpread{A constructed zero can pass weaker tests}{PROJECT COUNTERMODEL / NO ACTUAL TAU ZERO ESTABLISHED}{../shared/figures/synthetic_model.pdf}{
For any bound $B$ and modulus $M$, choose a prime $q>B$ with $q\equiv-1\pmod{50M}$ and $\tau(q)\equiv0\pmod{10M}$. Change its prime parameter to zero and retain all other prime parameters.

Extend by Hecke recurrence and multiplicativity. The sequence $f$ matches $\tau$ through $B$ and modulo $10M$ at every index. It retains the size bound, parity, modulo-$5$ congruence and primitive-square theorem.

Yet $f(q)=0$.}{This shows insufficiency of those premises. It does not claim unconditionally that $f$ differs from $\tau$ or is nonmodular.}

\GuideSpread{The global gap is a quantifier gap}{LOGICAL DISTINCTION / THE UNRESOLVED PRIME OBLIGATION}{../shared/figures/quantifiers.pdf}{
The proven statement chooses a prime after a modulus. That prime can change when the modulus changes.

Lehmer concerns each actual coefficient at its own prime. Divisibility by every integer at the same prime would force that coefficient to be zero.

Adaptive moduli and additional global consequences of $\Delta$ remain possible. The current routes do not supply the required uniform nonzero conclusion.}{The next useful step is a specific property of actual $\Delta$ excluding a zero prime parameter without presupposing Lehmer.}

\clearpage
{\fontsize{22}{26}\selectfont\bfseries\color{Ink}Sources and the evidence boundary}\par
{\sffamily\small\color{Gray}READ WITH THE COMPANION MONOGRAPH AND ITS PINNED REPORTS}\par
\vspace{3mm}
\begin{minipage}[t]{0.48\linewidth}\vspace{0pt}\small
\textbf{Classical source leaves.}\par
Han, \emph{The Nekrasov--Okounkov hook length formula}, arXiv:0805.1398v1, Theorems 1.1 and 1.2. Five-core coding, normalization and signed packets.\par
Balakrishnan--Craig--Ono, \emph{Variations of Lehmer's Conjecture for Ramanujan's tau-function}, arXiv:2005.10345v1: equation (1.2), p.1, for modulo $5$; Theorem 2.2, p.4, for Hecke laws and bound; \S3, p.7, for parity.\par
Deligne, \emph{Formes modulaires et repr\'esentations $\ell$-adiques}, Bourbaki expos\'e 355, 1968/69, published 1971, pp.139--172. Representation construction and Theorem 5.6.\par
Deligne, \emph{La conjecture de Weil I}, IH\'ES 43 (1974), Theorem 1.6. With the earlier implication, this yields the bound.\par
Serre, \emph{Congruences et formes modulaires}, Bourbaki expos\'e 416, 1971/72, published 1973, \S3.1: compatible representations, full-level unramified scope.\par
Garvan--Kim--Stanton, \emph{Cranks and $t$-cores}, Invent. Math. 101 (1990), author manuscript, Theorem 4 and equation (5.3), p.8: five-core count.\par
Stevenhagen--Lenstra, \emph{Chebotar\"ev and his density theorem}, 1996, pp.31,34: natural density. Colmez--Fontaine, \emph{Construction des repr\'esentations $p$-adiques semi-stables}, 2000, author manuscript \S3.4: weak admissibility.\par
Ibukiyama--Williams, arXiv:2505.23497v1, \S3.4, pp.18--19: harmonic theta modularity. Stein, \emph{Modular Forms of Level 1}, Theorem 2.11: valence. The vector-to-scalar descent is a project argument.\par
Exact URLs, versions and access limits are in the monograph bibliography and checked source contracts.
\end{minipage}\hfill
\begin{minipage}[t]{0.48\linewidth}\vspace{0pt}\small
\textbf{Project arguments and exact data.}\par
The 26 September packet, tree, content-transport and final-push archives are historical inputs. Their argument, computation, literature and definition reviews are distinct. Project status implies neither priority nor formal proof-assistant certification.\par
The coefficient table is the current independent integer replay through $512$. Shell figures are regenerated by integer enumeration. Logarithms and normalized heights convert data for display.\par
Threshold, primitive-square, content-conservation, interpolation and countermodel spreads summarize written arguments. Finite illustrations do not prove universal statements.\par
The trace-zero spread constructs a local filtered module. The fixed-modulus spread uses actual finite Galois images and Chebotarev. These two kinds of local information differ.\par
\textbf{Release status.} This program has not proved prime nonvanishing and has found no actual prime zero. The companion reader preserves the exact remaining implication.
\end{minipage}
\vfill
\colorbox{Soft}{\parbox{\dimexpr\linewidth-2\fboxsep\relax}{\small\textbf{Carry forward.} Exact examples, scoped theorems and route obstructions sharpen the problem while preserving its universal prime quantifier.}}
\end{document}
